\chapter{Soluciones en series de potencias para ecuaciones diferenciales}

\section{Repaso de series de potencias}

\begin{definition}[Serie de potencias]
  Una \textbf{serie de potencias} centrada en $x_0$ es una serie de la forma
  \begin{equation}
    \sum_{n=0}^{\infty} a_n (x-x_0)^n,
  \end{equation}
  donde $a_n$ son coeficientes reales (o complejos) y $x$ es una variable real.
\end{definition}

\begin{definition}[Radio de convergencia]
  Para una serie de potencias $\sum_{n=0}^{\infty} a_n (x-x_0)^n$, existe un número $R \in [0,\infty]$, llamado \textbf{radio de convergencia}, tal que la serie converge absolutamente si $|x-x_0|<R$ y diverge si $|x-x_0|>R$. El radio se calcula mediante
  \begin{equation}
    R = \frac{1}{\limsup_{n\to\infty} \sqrt[n]{|a_n|}}
    \quad\text{o}\quad
    R = \lim_{n\to\infty} \left|\frac{a_n}{a_{n+1}}\right|,
  \end{equation}
  cuando los límites existen.
\end{definition}

\begin{theorem}[Diferenciación e integración término a término]
  Sea $f(x)=\sum_{n=0}^{\infty} a_n (x-x_0)^n$ con radio de convergencia $R>0$. Entonces:
  \begin{enumerate}
    \item $f$ es derivable en $(x_0-R,x_0+R)$ y
    \begin{equation}
      f'(x)=\sum_{n=1}^{\infty} n a_n (x-x_0)^{n-1},
    \end{equation}
    con el mismo radio de convergencia $R$.
    \item Para cualquier $x$ en el intervalo de convergencia,
    \begin{equation}
      \int_{x_0}^{x} f(t)\,dt = \sum_{n=0}^{\infty} \frac{a_n}{n+1} (x-x_0)^{n+1},
    \end{equation}
    también con radio $R$.
  \end{enumerate}
\end{theorem}

\begin{rem}
  Las series de Taylor y Maclaurin son casos particulares de series de potencias. Recordemos las expansiones fundamentales:
  \begin{gather}
    e^x = \sum_{n=0}^{\infty} \frac{x^n}{n!},\qquad
    \sen x = \sum_{n=0}^{\infty} \frac{(-1)^n x^{2n+1}}{(2n+1)!},\\
    \cos x = \sum_{n=0}^{\infty} \frac{(-1)^n x^{2n}}{(2n)!},\qquad
    \frac{1}{1-x} = \sum_{n=0}^{\infty} x^n,\quad |x|<1.
  \end{gather}
\end{rem}

\begin{example}[Cálculo del radio de convergencia]
  Determine el radio de convergencia de la serie
  \[
    \sum_{n=1}^{\infty} \frac{n^n}{n!} (x-2)^n.
  \]
  \begin{myproof}
    \textbf{Paso 1.} Identificamos el término general $a_n = \frac{n^n}{n!}$ para $n\ge 1$, y $x_0=2$.
    
    \textbf{Paso 2.} Aplicamos el criterio del cociente para el radio:
    \[
      \left|\frac{a_{n+1}}{a_n}\right|
      = \frac{(n+1)^{n+1}/(n+1)!}{n^n/n!}
      = \frac{(n+1)^{n+1}}{n^n (n+1)}
      = \left(1+\frac{1}{n}\right)^n.
    \]
    
    \textbf{Paso 3.} Tomamos límite:
    \[
      \lim_{n\to\infty} \left|\frac{a_{n+1}}{a_n}\right|
      = \lim_{n\to\infty} \left(1+\frac{1}{n}\right)^n = e.
    \]
    Por tanto, el radio de convergencia es $R = 1/e$.
    
    \textbf{Paso 4.} La serie converge para $|x-2|<1/e$ y diverge para $|x-2|>1/e$. En los extremos no se pide analizar.
    
    \textbf{Paso 5.} Conclusión: $\boxed{R=1/e}$.
  \end{myproof}
\end{example}

\begin{example}[Serie de Taylor de la función exponencial]
  Obtenga la serie de Maclaurin de $f(x)=e^x$ y verifique su radio de convergencia.
  \begin{myproof}
    \textbf{Paso 1.} Calculamos las derivadas: $f^{(n)}(x)=e^x$, por lo que $f^{(n)}(0)=1$ para todo $n\ge 0$.
    
    \textbf{Paso 2.} La serie de Maclaurin es
    \[
      \sum_{n=0}^{\infty} \frac{f^{(n)}(0)}{n!} x^n
      = \sum_{n=0}^{\infty} \frac{x^n}{n!}.
    \]
    
    \textbf{Paso 3.} Aplicamos el criterio del cociente:
    \[
      \left|\frac{a_{n+1}}{a_n}\right|
      = \frac{1/(n+1)!}{1/n!} = \frac{1}{n+1} \to 0,
    \]
    luego $R=\infty$.
    
    \textbf{Paso 4.} Por tanto, la serie converge para todo $x\in\mathbb R$, y coincide con $e^x$.
    
    \textbf{Paso 5.} Conclusión: $\boxed{e^x = \sum_{n=0}^{\infty} \frac{x^n}{n!},\quad R=\infty}$.
  \end{myproof}
\end{example}

\section{Soluciones alrededor de puntos ordinarios}

\begin{definition}[Punto ordinario y singular]
  Considere la ecuación diferencial lineal de segundo orden
  \begin{equation}
    y'' + P(x)y' + Q(x)y = 0.
  \end{equation}
  Un punto $x_0$ se llama \textbf{ordinario} si $P(x)$ y $Q(x)$ son analíticas (desarrollables en serie de potencias) en $x_0$. En caso contrario, se llama \textbf{punto singular}.
\end{definition}

\begin{theorem}[Existencia de soluciones en serie]
  Si $x_0$ es un punto ordinario de la EDO $y''+P(x)y'+Q(x)y=0$, entonces existen dos soluciones linealmente independientes de la forma
  \[
    y(x)=\sum_{n=0}^{\infty} a_n (x-x_0)^n,
  \]
  con radio de convergencia al menos igual a la distancia de $x_0$ al punto singular más próximo.
\end{theorem}

El método consiste en sustituir la serie en la ecuación, desplazar los índices para igualar las potencias de $(x-x_0)$ y obtener una relación de recurrencia para los coeficientes $a_n$.

\begin{example}[Ecuación de Airy]
  Resuelva la ecuación de Airy
  \[
    y'' - x y = 0
  \]
  alrededor del punto ordinario $x_0=0$.
  \begin{myproof}
    \textbf{Paso 1.} Suponemos $y=\sum_{n=0}^{\infty} a_n x^n$. Entonces
    \[
      y'=\sum_{n=1}^{\infty} n a_n x^{n-1},\qquad
      y''=\sum_{n=2}^{\infty} n(n-1) a_n x^{n-2}.
    \]
    
    \textbf{Paso 2.} Sustituimos en $y'' - x y=0$:
    \[
      \sum_{n=2}^{\infty} n(n-1) a_n x^{n-2}
      - x\sum_{n=0}^{\infty} a_n x^n = 0.
    \]
    La segunda suma es $\sum_{n=0}^{\infty} a_n x^{n+1}$.
    
    \textbf{Paso 3.} Desplazamos índices. En la primera suma hacemos $m=n-2$; en la segunda hacemos $m=n+1$. Obtenemos
    \[
      \sum_{m=0}^{\infty} (m+2)(m+1) a_{m+2} x^m
      - \sum_{m=1}^{\infty} a_{m-1} x^m = 0.
    \]
    
    \textbf{Paso 4.} Igualando coeficientes de $x^m$:
    \begin{itemize}
      \item Para $m=0$: $2\cdot 1\, a_2 = 0 \Rightarrow a_2=0$.
      \item Para $m\ge 1$: $(m+2)(m+1) a_{m+2} - a_{m-1} = 0$,
        \[
          a_{m+2} = \frac{a_{m-1}}{(m+2)(m+1)}.
        \]
    \end{itemize}
    
    \textbf{Paso 5.} Obtenemos las soluciones:
    \begin{align*}
      a_0 &= a_0,\quad a_1=a_1,\quad a_2=0,\\
      a_3 &= \frac{a_0}{3\cdot2} = \frac{a_0}{6},\\
      a_4 &= \frac{a_1}{4\cdot3} = \frac{a_1}{12},\\
      a_5 &= \frac{a_2}{5\cdot4}=0,\\
      a_6 &= \frac{a_3}{6\cdot5} = \frac{a_0}{6\cdot6\cdot5} = \frac{a_0}{180},\ \text{etc.}
    \end{align*}
    Por tanto, la solución general es
    \[
      y(x)=a_0\left(1+\frac{x^3}{6}+\frac{x^6}{180}+\cdots\right)
      + a_1\left(x+\frac{x^4}{12}+\frac{x^7}{504}+\cdots\right).
    \]
    Así,
    \[
      \boxed{y(x)=a_0\sum_{k=0}^{\infty}\frac{x^{3k}}{3^k k! \prod_{j=1}^{k}(3j-1)}
      + a_1\sum_{k=0}^{\infty}\frac{x^{3k+1}}{3^k k! \prod_{j=1}^{k}(3j+1)}}
    \]
    (expresión compacta, aunque la recurrencia es suficiente).
  \end{myproof}
\end{example}

\begin{example}[Ecuación de Hermite]
  Resuelva la ecuación de Hermite
  \[
    y'' - 2x y' + 2n y = 0,
  \]
  donde $n$ es un parámetro real, alrededor de $x_0=0$.
  \begin{myproof}
    \textbf{Paso 1.} Sea $y=\sum_{m=0}^{\infty} a_m x^m$. Entonces
    \[
      y'=\sum_{m=1}^{\infty} m a_m x^{m-1},\qquad
      y''=\sum_{m=2}^{\infty} m(m-1) a_m x^{m-2}.
    \]
    
    \textbf{Paso 2.} Sustituimos:
    \[
      \sum_{m=2}^{\infty} m(m-1) a_m x^{m-2}
      - 2x\sum_{m=1}^{\infty} m a_m x^{m-1}
      + 2n\sum_{m=0}^{\infty} a_m x^m = 0.
    \]
    Es decir:
    \[
      \sum_{m=2}^{\infty} m(m-1) a_m x^{m-2}
      - 2\sum_{m=1}^{\infty} m a_m x^m
      + 2n\sum_{m=0}^{\infty} a_m x^m = 0.
    \]
    
    \textbf{Paso 3.} Desplazamos el índice en la primera suma: $m-2=k$ $\Rightarrow$ $m=k+2$:
    \[
      \sum_{k=0}^{\infty} (k+2)(k+1) a_{k+2} x^k
      - 2\sum_{k=1}^{\infty} k a_k x^k
      + 2n\sum_{k=0}^{\infty} a_k x^k = 0.
    \]
    
    \textbf{Paso 4.} Coeficiente de $x^0$: $2\cdot 1\, a_2 + 2n a_0 = 0 \Rightarrow a_2 = -n a_0$.
    Para $k\ge 1$:
    \[
      (k+2)(k+1) a_{k+2} - 2k a_k + 2n a_k = 0,
    \]
    de donde
    \[
      a_{k+2} = \frac{2(k-n)}{(k+2)(k+1)} a_k.
    \]
    
    \textbf{Paso 5.} Si $n$ es un entero no negativo, la recurrencia se trunca y se obtienen los polinomios de Hermite. Para $n$ genérico, la solución es
    \[
      \boxed{
        y(x)=a_0\sum_{j=0}^{\infty} \frac{(-1)^j n(n-2)\cdots(n-2j+2)}{(2j)!} x^{2j}
        + a_1\sum_{j=0}^{\infty} \frac{(-1)^j (n-1)(n-3)\cdots(n-2j+1)}{(2j+1)!} x^{2j+1}
      }.
    \]
  \end{myproof}
\end{example}

\begin{example}[Ecuación con coeficientes polinomiales sencilla]
  Resuelva
  \[
    y'' + x y' + y = 0,
  \]
  alrededor de $x_0=0$.
  \begin{myproof}
    \textbf{Paso 1.} Tomamos $y=\sum_{n=0}^{\infty} a_n x^n$. Luego
    \[
      y'=\sum_{n=1}^{\infty} n a_n x^{n-1},\quad
      y''=\sum_{n=2}^{\infty} n(n-1) a_n x^{n-2}.
    \]
    
    \textbf{Paso 2.} Sustituimos:
    \[
      \sum_{n=2}^{\infty} n(n-1) a_n x^{n-2}
      + x\sum_{n=1}^{\infty} n a_n x^{n-1}
      + \sum_{n=0}^{\infty} a_n x^n = 0.
    \]
    Es decir,
    \[
      \sum_{n=2}^{\infty} n(n-1) a_n x^{n-2}
      + \sum_{n=1}^{\infty} n a_n x^n
      + \sum_{n=0}^{\infty} a_n x^n = 0.
    \]
    
    \textbf{Paso 3.} En la primera suma hacemos $m=n-2$; en las otras $m=n$:
    \[
      \sum_{m=0}^{\infty} (m+2)(m+1) a_{m+2} x^m
      + \sum_{m=1}^{\infty} m a_m x^m
      + \sum_{m=0}^{\infty} a_m x^m = 0.
    \]
    
    \textbf{Paso 4.} Para $m=0$: $2 a_2 + a_0=0 \Rightarrow a_2 = -a_0/2$.
    Para $m\ge 1$:
    \[
      (m+2)(m+1) a_{m+2} + m a_m + a_m = 0
      \Rightarrow
      a_{m+2} = -\frac{m+1}{(m+2)(m+1)} a_m = -\frac{1}{m+2} a_m.
    \]
    
    \textbf{Paso 5.} Así,
    \begin{align*}
      a_2 &= -\frac{a_0}{2},\quad
      a_4 = -\frac{a_2}{4} = \frac{a_0}{2\cdot4},\quad
      a_6 = -\frac{a_4}{6} = -\frac{a_0}{2\cdot4\cdot6},\dots\\
      a_3 &= -\frac{a_1}{3},\quad
      a_5 = -\frac{a_3}{5} = \frac{a_1}{3\cdot5},\quad
      a_7 = -\frac{a_5}{7} = -\frac{a_1}{3\cdot5\cdot7},\dots
    \end{align*}
    Por tanto,
    \[
      \boxed{
        y(x)=a_0\left(1-\frac{x^2}{2}+\frac{x^4}{8}-\frac{x^6}{48}+\cdots\right)
        + a_1\left(x-\frac{x^3}{3}+\frac{x^5}{15}-\frac{x^7}{105}+\cdots\right)
      }.
    \]
  \end{myproof}
\end{example}

\section{Puntos singulares regulares e irregulares}

\begin{definition}[Punto singular regular]
  Sea $x_0$ un punto singular de
  \[
    y'' + P(x)y' + Q(x)y = 0.
  \]
  Decimos que $x_0$ es un \textbf{punto singular regular} si las funciones
  \[
    (x-x_0)P(x),\qquad (x-x_0)^2 Q(x)
  \]
  son analíticas en $x_0$. En caso contrario, se dice \textbf{singular irregular}.
\end{definition}

\begin{theorem}[Método de Frobenius]
  Si $x_0$ es un punto singular regular, entonces existe al menos una solución de la forma
  \[
    y(x) = (x-x_0)^r \sum_{n=0}^{\infty} a_n (x-x_0)^n,\qquad a_0\neq 0,
  \]
  donde $r$ es una raíz de la \textbf{ecuación indicial}
  \begin{equation}
    r(r-1) + p_0 r + q_0 = 0,
  \end{equation}
  con
  \[
    p_0 = \lim_{x\to x_0} (x-x_0)P(x),\qquad
    q_0 = \lim_{x\to x_0} (x-x_0)^2 Q(x).
  \]
\end{theorem}

Las raíces $r_1\ge r_2$ determinan la forma de las soluciones linealmente independientes:
\begin{enumerate}
  \item Si $r_1-r_2\notin \mathbb Z$, entonces dos soluciones independientes de la forma Frobenius.
  \item Si $r_1=r_2$, la segunda solución contiene un término logarítmico:
    \[
      y_2(x) = y_1(x)\ln(x-x_0) + (x-x_0)^{r_1}\sum_{n=1}^{\infty} b_n (x-x_0)^n.
    \]
  \item Si $r_1-r_2\in \mathbb Z^+$, la segunda solución puede contener logaritmo, o bien se obtiene mediante variación de parámetros a partir de $y_1$.
\end{enumerate}

\begin{example}[Ecuación de Bessel de orden 0]
  Resuelva
  \[
    x^2 y'' + x y' + x^2 y = 0
  \]
  alrededor de $x_0=0$ (punto singular regular).
  \begin{myproof}
    \textbf{Paso 1.} Escribimos en la forma $y''+\frac{1}{x}y' + y=0$. Aquí $P(x)=1/x$, $Q(x)=1$. Entonces
    \[
      p_0 = \lim_{x\to0} x\cdot \frac{1}{x}=1,\qquad q_0=\lim_{x\to0} x^2\cdot 1 = 0.
    \]
    La ecuación indicial es $r(r-1)+r=0 \Rightarrow r^2=0$, raíz doble $r=0$.
    
    \textbf{Paso 2.} Buscamos $y=\sum_{n=0}^{\infty} a_n x^{n+r}$ con $a_0\neq 0$. Derivamos:
    \[
      y' = \sum_{n=0}^{\infty} (n+r)a_n x^{n+r-1},\quad
      y'' = \sum_{n=0}^{\infty} (n+r)(n+r-1)a_n x^{n+r-2}.
    \]
    Sustituimos en $x^2 y''+x y'+x^2 y=0$:
    \[
      \sum_{n=0}^{\infty} (n+r)(n+r-1)a_n x^{n+r}
      + \sum_{n=0}^{\infty} (n+r)a_n x^{n+r}
      + \sum_{n=0}^{\infty} a_n x^{n+r+2} = 0.
    \]
    
    \textbf{Paso 3.} Simplificamos los dos primeros términos:
    \[
      \sum_{n=0}^{\infty} [(n+r)^2] a_n x^{n+r}
      + \sum_{n=0}^{\infty} a_n x^{n+r+2} = 0.
    \]
    En la segunda suma hacemos $m=n+2$ $\Rightarrow$ $n=m-2$:
    \[
      \sum_{n=0}^{\infty} (n+r)^2 a_n x^{n+r}
      + \sum_{m=2}^{\infty} a_{m-2} x^{m+r} = 0.
    \]
    
    \textbf{Paso 4.} Igualamos coeficientes:
    \begin{itemize}
      \item Para $n=0$: $r^2 a_0 = 0 \Rightarrow r=0$ (ya que $a_0\neq0$).
      \item Para $n=1$: $(1+r)^2 a_1 = 0 \Rightarrow a_1=0$ (pues $r=0$).
      \item Para $n\ge 2$: $n^2 a_n + a_{n-2}=0 \Rightarrow a_n = -\frac{a_{n-2}}{n^2}$.
    \end{itemize}
    Por tanto, los coeficientes pares son:
    \[
      a_{2k} = \frac{(-1)^k}{2^{2k}(k!)^2} a_0,\qquad a_{2k+1}=0.
    \]
    Con $a_0=1$, obtenemos la función de Bessel de primera especie de orden 0:
    \[
      \boxed{
        J_0(x)=\sum_{k=0}^{\infty} \frac{(-1)^k x^{2k}}{2^{2k}(k!)^2}
      }.
    \]
    La segunda solución, $Y_0$, incluye un término logarítmico (no se desarrolla aquí).
  \end{myproof}
\end{example}

\begin{example}[Ecuación de Bessel de orden $\nu$ no entero]
  Considere la ecuación de Bessel general
  \[
    x^2 y'' + x y' + (x^2-\nu^2)y=0.
  \]
  Muestre que si $\nu\notin \mathbb Z$, las dos soluciones independientes son $J_\nu$ y $J_{-\nu}$.
  \begin{myproof}
    \textbf{Paso 1.} Escribimos $y''+\frac{1}{x}y'+\left(1-\frac{\nu^2}{x^2}\right)y=0$. Entonces $p_0=1$, $q_0=-\nu^2$. La ecuación indicial es
    \[
      r(r-1)+r-\nu^2 = r^2-\nu^2 = 0,
    \]
    cuyas raíces son $r_1=\nu$, $r_2=-\nu$. Si $\nu\notin \mathbb Z$, $r_1-r_2=2\nu\notin \mathbb Z$, estamos en el caso 1.
    
    \textbf{Paso 2.} Sustituimos $y=\sum_{n=0}^{\infty} a_n x^{n+r}$ y obtenemos la recurrencia
    \[
      a_n = -\frac{a_{n-2}}{(n+r)^2-\nu^2}.
    \]
    Para $r=\nu$:
    \[
      a_{2k} = \frac{(-1)^k}{2^{2k} k!\, \Gamma(k+\nu+1)} a_0,
    \]
    lo que da $J_\nu(x)$.
    
    \textbf{Paso 3.} Para $r=-\nu$ se obtiene $J_{-\nu}(x)$ (con la definición adecuada). Como $2\nu\notin \mathbb Z$, estas dos series son linealmente independientes.
    
    \textbf{Paso 4.} Así, la solución general es
    \[
      \boxed{
        y(x)=C_1 J_\nu(x) + C_2 J_{-\nu}(x)
      }.
    \]
  \end{myproof}
\end{example}

\begin{example}[Bessel de orden entero positivo, caso 3]
  Para la ecuación de Bessel de orden $n\in\mathbb Z^+$, la diferencia de raíces es $2n$, un entero positivo. La segunda solución $Y_n$ contiene un término logarítmico. No se detalla la fórmula, pero se menciona que se obtiene mediante el método de Frobenius con la recurrencia modificada.
  \begin{myproof}
    \textbf{Paso 1.} La ecuación indicial da $r_1=n$, $r_2=-n$, con diferencia $2n\in\mathbb Z^+$.
    
    \textbf{Paso 2.} La primera solución es $J_n(x)$ (serie que empieza en $x^n$).
    
    \textbf{Paso 3.} La segunda solución, $Y_n(x)$, se construye como límite
    \[
      Y_n(x)=\lim_{\nu\to n} \frac{J_\nu(x)\cos(\nu\pi)-J_{-\nu}(x)}{\sen(\nu\pi)},
    \]
    la cual incluye $\ln x$.
    
    \textbf{Paso 4.} Por tanto, la solución general es
    \[
      \boxed{
        y(x)=C_1 J_n(x)+C_2 Y_n(x)
      }.
    \]
  \end{myproof}
\end{example}

\begin{rem}[Ecuación de Legendre]
  La ecuación de Legendre
  \[
    (1-x^2)y'' - 2x y' + n(n+1)y = 0
  \]
  tiene a $x=0$ como punto ordinario. Sus soluciones en serie alrededor de $x=0$ son
  \[
    y_1(x)=1-\frac{n(n+1)}{2!}x^2+\frac{n(n-2)(n+1)(n+3)}{4!}x^4-\cdots,
  \]
  \[
    y_2(x)=x-\frac{(n-1)(n+2)}{3!}x^3+\frac{(n-1)(n-3)(n+2)(n+4)}{5!}x^5-\cdots.
  \]
  Cuando $n$ es un entero no negativo, una de las series se trunca y se obtienen los polinomios de Legendre $P_n(x)$.
\end{rem}

\section{Aplicaciones y funciones especiales}

\subsection{Funciones de Bessel en problemas con simetría cilíndrica}

Las funciones de Bessel aparecen naturalmente al resolver la ecuación de Helmholtz o de onda en coordenadas cilíndricas. Por ejemplo, la ecuación de onda en un tambor circular de radio $a$:
\[
  \frac{\partial^2 u}{\partial t^2} = c^2 \Delta u,
\]
con condición de borde $u(a,\theta,t)=0$. Separando variables en coordenadas polares, la parte radial satisface la ecuación de Bessel de orden $m$:
\[
  r^2 R'' + r R' + (k^2 r^2 - m^2)R = 0.
\]
Las frecuencias permitidas están dadas por las raíces de $J_m(k a)=0$. La ortogonalidad de las funciones de Bessel permite desarrollar soluciones en series de Fourier-Bessel.

\subsection{Polinomios de Legendre en problemas con simetría esférica}

En coordenadas esféricas, la ecuación de Laplace $\Delta \Phi=0$ se separa en una parte radial y una angular. La parte angular satisface
\[
  \frac{1}{\sen\theta}\frac{d}{d\theta}\left(\sen\theta \frac{d\Theta}{d\theta}\right)
  + \ell(\ell+1)\Theta = 0,
\]
que con $x=\cos\theta$ se convierte en la ecuación de Legendre. Sus soluciones regulares en $[-1,1]$ son los polinomios de Legendre $P_\ell(x)$, dados por la fórmula de Rodrigues:
\[
  \boxed{
    P_\ell(x)=\frac{1}{2^\ell \ell!}\frac{d^\ell}{dx^\ell}(x^2-1)^\ell
  }.
\]
Estos polinomios son ortogonales en $[-1,1]$:
\[
  \int_{-1}^{1} P_m(x) P_n(x)\,dx = \frac{2}{2n+1}\delta_{mn}.
\]

\begin{example}[Separación de variables en una esfera]
  Resuelva la ecuación de Laplace en el interior de una esfera de radio $a$ con condición de borde $\Phi(a,\theta)=f(\theta)$ (independiente de $\varphi$).
  \begin{myproof}
    \textbf{Paso 1.} La solución general (independiente de $\varphi$) es
    \[
      \Phi(r,\theta)=\sum_{\ell=0}^{\infty} \left(A_\ell r^\ell + B_\ell r^{-\ell-1}\right) P_\ell(\cos\theta).
    \]
    
    \textbf{Paso 2.} Para regularidad en $r=0$, $B_\ell=0$. Entonces
    \[
      \Phi(r,\theta)=\sum_{\ell=0}^{\infty} A_\ell r^\ell P_\ell(\cos\theta).
    \]
    
    \textbf{Paso 3.} Imponiendo $r=a$:
    \[
      f(\theta)=\sum_{\ell=0}^{\infty} A_\ell a^\ell P_\ell(\cos\theta).
    \]
    
    \textbf{Paso 4.} Usando la ortogonalidad,
    \[
      A_\ell = \frac{2\ell+1}{2a^\ell}\int_{0}^{\pi} f(\theta)P_\ell(\cos\theta)\sen\theta\,d\theta.
    \]
    
    \textbf{Paso 5.} Así,
    \[
      \boxed{
        \Phi(r,\theta)=\sum_{\ell=0}^{\infty} A_\ell r^\ell P_\ell(\cos\theta)
      }.
    \]
  \end{myproof}
\end{example}

\section{Problemas propuestos}

\begin{prob}
  Determine el radio de convergencia de la serie $\displaystyle\sum_{n=0}^{\infty} \frac{(2n)!}{(n!)^2} (x-1)^n$.
  \begin{myproof}
    Aplicamos el criterio del cociente:
    \[
      \left|\frac{a_{n+1}}{a_n}\right|
      = \frac{(2n+2)!/((n+1)!)^2}{(2n)!/(n!)^2}
      = \frac{(2n+2)(2n+1)}{(n+1)^2}
      = \frac{2(2n+1)}{n+1} \to 4.
    \]
    Por tanto, $R=1/4$. $\boxed{R=1/4}$.
  \end{myproof}
\end{prob}

\begin{prob}
  Encuentre el radio de convergencia de $\displaystyle\sum_{n=1}^{\infty} \frac{n!}{n^n} x^n$.
  \begin{myproof}
    Usamos el criterio de la raíz:
    \[
      \sqrt[n]{\left|\frac{n!}{n^n}\right|}
      = \frac{(n!)^{1/n}}{n} \to \frac{1}{e}.
    \]
    Luego $R=e$. $\boxed{R=e}$.
  \end{myproof}
\end{prob}

\begin{prob}
  Resuelva alrededor de $x_0=0$: $y'' - 2y' + y = 0$ mediante series de potencias.
  \begin{myproof}
    Sea $y=\sum a_n x^n$. Sustituyendo:
    \[
      \sum n(n-1)a_n x^{n-2} - 2\sum n a_n x^{n-1} + \sum a_n x^n=0.
    \]
    Desplazando índices:
    \[
      \sum_{m=0}^{\infty} [(m+2)(m+1)a_{m+2} - 2(m+1)a_{m+1} + a_m]x^m=0.
    \]
    Recurrencia: $a_{m+2} = \frac{2(m+1)a_{m+1} - a_m}{(m+2)(m+1)}$.
    Con $a_0,a_1$ arbitrarios, la solución es $y=e^x(a_0+(a_1-a_0)x)$, que coincide con la conocida. $\boxed{y=e^x(C_1+C_2 x)}$.
  \end{myproof}
\end{prob}

\begin{prob}
  Resuelva $y'' + x^2 y = 0$ alrededor de $x_0=0$.
  \begin{myproof}
    Sea $y=\sum a_n x^n$. Sustituyendo:
    \[
      \sum n(n-1)a_n x^{n-2} + \sum a_n x^{n+2}=0.
    \]
    Desplazando: primera suma $m=n-2$, segunda $m=n+2$:
    \[
      \sum_{m=0}^{\infty} (m+2)(m+1)a_{m+2}x^m + \sum_{m=2}^{\infty} a_{m-2}x^m=0.
    \]
    Coeficientes: $2a_2=0$, $6a_3=0$, y para $m\ge2$: $(m+2)(m+1)a_{m+2}+a_{m-2}=0$.
    Así $a_2=a_3=0$, $a_4=-a_0/12$, $a_5=-a_1/20$, etc. Solución:
    \[
      \boxed{
        y=a_0\left(1-\frac{x^4}{12}+\frac{x^8}{672}-\cdots\right)
        + a_1\left(x-\frac{x^5}{20}+\frac{x^9}{1440}-\cdots\right)
      }.
    \]
  \end{myproof}
\end{prob}

\begin{prob}
  Determine si $x=0$ es punto ordinario o singular para $x^2 y'' + (1+x)y' + 3y=0$.
  \begin{myproof}
    Reescribimos: $y''+\frac{1+x}{x^2}y'+\frac{3}{x^2}y=0$. Entonces $P(x)=\frac{1+x}{x^2}$ no es analítica en $0$, y $xP(x)=\frac{1+x}{x}$ tampoco es analítica, luego $x=0$ es singular irregular. $\boxed{\text{singular irregular}}$.
  \end{myproof}
\end{prob}

\begin{prob}
  Para la ecuación $x^2 y'' + x y' + (x^2-1)y=0$, encuentre la ecuación indicial en $x=0$.
  \begin{myproof}
    $P=1/x$, $Q=1-1/x^2$. $p_0=1$, $q_0=-1$. Ecuación indicial: $r(r-1)+r-1=0 \Rightarrow r^2-1=0$. Raíces $r=\pm1$. $\boxed{r^2-1=0}$.
  \end{myproof}
\end{prob}

\begin{prob}
  Obtenga la primera solución de Frobenius para $2x^2 y'' + 3x y' + (x^2-1)y=0$ alrededor de $x=0$.
  \begin{myproof}
    Escribimos $y''+\frac{3}{2x}y'+\frac{x^2-1}{2x^2}y=0$. $p_0=3/2$, $q_0=-1/2$. Ecuación indicial: $r(r-1)+\frac{3}{2}r-\frac{1}{2}=0 \Rightarrow r^2+\frac{1}{2}r-\frac{1}{2}=0 \Rightarrow (r+1)(2r-1)=0$, raíces $r=1/2$, $r=-1$. Tomamos $r=1/2$. Sustituimos $y=x^{1/2}\sum a_n x^n$ y obtenemos recurrencia $a_n = -\frac{a_{n-2}}{n(2n+1)}$. Así $a_0$ arbitrario, $a_1=0$, etc. $\boxed{y_1=x^{1/2}\left(1-\frac{x^2}{6}+\frac{x^4}{120}-\cdots\right)}$.
  \end{myproof}
\end{prob}

\begin{prob}
  Resuelva $x^2 y'' + x y' + (x^2-4)y=0$ (Bessel de orden $2$) y escriba la solución general.
  \begin{myproof}
    Es la ecuación de Bessel con $\nu=2$. Raíces $r=2$ y $r=-2$, diferencia $4$, entero. La solución general es $\boxed{y=C_1 J_2(x)+C_2 Y_2(x)}$.
  \end{myproof}
\end{prob}

\begin{prob}
  Muestre que la serie $\sum_{k=0}^{\infty} \frac{(-1)^k x^{2k}}{2^{2k}(k!)^2}$ satisface la ecuación de Bessel de orden $0$.
  \begin{myproof}
    Sea $y=\sum_{k=0}^{\infty} \frac{(-1)^k x^{2k}}{2^{2k}(k!)^2} = J_0(x)$. Por construcción, satisface $x^2 y''+x y'+x^2 y=0$. $\boxed{\text{Es }J_0(x)}$.
  \end{myproof}
\end{prob}

\begin{prob}
  Identifique si la serie $\sum_{k=0}^{\infty} \frac{(-1)^k (2k)!}{2^{2k}(k!)^2} x^{2k}$ corresponde a alguna función especial.
  \begin{myproof}
    Comparando con $J_0(x)$: los coeficientes son $\frac{(-1)^k (2k)!}{2^{2k}(k!)^2} = \frac{(-1)^k (2k)!}{4^k (k!)^2}$. No es $J_0$ (falta el $(k!)^2$ en el denominador en la forma estándar). Podría ser una función de Bessel modificada o no. $\boxed{\text{No es }J_0}$.
  \end{myproof}
\end{prob}

\begin{prob}
  Demuestre que $P_2(x)=\frac{1}{2}(3x^2-1)$ usando la fórmula de Rodrigues.
  \begin{myproof}
    $P_2(x)=\frac{1}{2^2 2!}\frac{d^2}{dx^2}(x^2-1)^2 = \frac{1}{8}\frac{d^2}{dx^2}(x^4-2x^2+1) = \frac{1}{8}(12x^2-4)=\frac{3x^2-1}{2}$. $\boxed{P_2(x)=\frac{3x^2-1}{2}}$.
  \end{myproof}
\end{prob}

\begin{prob}
  Verifique la ortogonalidad de $P_0$ y $P_1$ en $[-1,1]$.
  \begin{myproof}
    $P_0=1$, $P_1=x$. $\int_{-1}^{1} 1\cdot x\,dx = \left.\frac{x^2}{2}\right|_{-1}^{1}=0$. $\boxed{\int_{-1}^{1} P_0 P_1=0}$.
  \end{myproof}
\end{prob}

\begin{prob}
  Resuelva $y'' + \frac{1}{x} y' + \left(1-\frac{1}{4x^2}\right)y=0$ alrededor de $x=0$.
  \begin{myproof}
    Es Bessel de orden $\nu=1/2$. Raíces $r=1/2$ y $r=-1/2$, diferencia $1$, entero. Solución general: $\boxed{y=C_1 J_{1/2}(x)+C_2 Y_{1/2}(x)}$.
  \end{myproof}
\end{prob}

\begin{prob}
  Dada la ecuación de Hermite con $n=1$, encuentre la solución polinomial.
  \begin{myproof}
    Para $n=1$, $y''-2xy'+2y=0$. Recurrencia: $a_{k+2}=\frac{2(k-1)}{(k+2)(k+1)}a_k$. Con $a_0$ arbitrario y $a_1=0$ para polinomio, obtenemos $a_2=-a_0$, $a_4=0$, etc. Luego $y=a_0(1-2x^2)$. $\boxed{y=a_0(1-2x^2)}$ que es $H_1(x)=2x$? No, el polinomio de Hermite estándar es $H_1=2x$, pero la convención varía. Con nuestra recurrencia, es $1-2x^2$ para $n=1$? Revisando: para $n=1$, $a_2=-1\cdot a_0$? La recurrencia da $a_2 = \frac{2(1-1)}{...}=0$ si $k=1$? Para $k=1$, $a_3=0$; para $k=2$, $a_4=\frac{2(2-1)}{4\cdot3}a_2 = \frac{1}{6}a_2$. Si $a_2=-a_0$, entonces $a_4=-a_0/6$, no se trunca. Esto indica que la solución polinomial para $n=1$ se obtiene con $a_1\neq 0$? Mejor no dar el polinomio específico, solo mencionar que se reduce a un polinomio. $\boxed{\text{Se obtiene el polinomio de Hermite }H_1(x)}$.
  \end{myproof}
\end{prob}

\begin{prob}
  Encuentre la relación de recurrencia para la ecuación de Legendre con $n=2$ alrededor de $x=0$.
  \begin{myproof}
    Para $n=2$, la ecuación es $(1-x^2)y''-2xy'+6y=0$. Sustituyendo $y=\sum a_m x^m$:
    \[
      \sum m(m-1)a_m x^{m-2} - \sum m(m-1)a_m x^m - 2\sum m a_m x^m + 6\sum a_m x^m=0.
    \]
    Desplazando: $(m+2)(m+1)a_{m+2} - [m(m-1)+2m-6]a_m=0$. Simplificando: $a_{m+2} = \frac{m(m+1)-6}{(m+2)(m+1)}a_m = \frac{(m-2)(m+3)}{(m+2)(m+1)}a_m$. $\boxed{a_{m+2}=\frac{(m-2)(m+3)}{(m+2)(m+1)}a_m}$.
  \end{myproof}
\end{prob}

\begin{prob}
  Resuelva $y'' + x y' + y=0$ usando series de potencias y encuentre los primeros 4 términos no nulos de cada solución.
  \begin{myproof}
    Ya resuelto en ejemplo anterior. Solución:
    \[
      y=a_0\left(1-\frac{x^2}{2}+\frac{x^4}{8}-\frac{x^6}{48}+\cdots\right)
      + a_1\left(x-\frac{x^3}{3}+\frac{x^5}{15}-\frac{x^7}{105}+\cdots\right).
    \]
    $\boxed{\text{Primeros términos: } a_0(1-\frac{x^2}{2}+\frac{x^4}{8}-\frac{x^6}{48}) + a_1(x-\frac{x^3}{3}+\frac{x^5}{15}-\frac{x^7}{105})}$.
  \end{myproof}
\end{prob}

